% ==============================================================================
% SECTION 2: SEARCH METHODOLOGY
% ==============================================================================
% Extensión: 300–500 palabras
% Estructura: pregunta PICO, bases de datos, ecuación de búsqueda,
% criterios de inclusión/exclusión, diagrama PRISMA (TikZ)

\section{Metodología de búsqueda}\label{sec:methods}

\subsection{Pregunta de investigación}

Se formuló una pregunta de investigación basada en el marco PICO:

¿Qué tan efectivos son los sistemas basados en Large Language Models (LLMs) y 
agentes en comparación con workflows tradicionales para la automatización de 
pipelines bioinformáticos?

\begin{table}[htbp]
  \centering
  \caption{Marco PICO de la pregunta de investigación.}\label{tab:pico}
  \begin{tabular}{@{}p{3.2cm} p{9.2cm}@{}}
    \toprule
    \textbf{Componente} & \textbf{Descripción} \\
    \midrule
    \textbf{P: Población}    & Pipelines bioinformáticos, incluyendo análisis de RNA-seq, genómica y transcriptómica. \\
    \textbf{I: Intervención} & Sistemas basados en Large Language Models (LLMs) y agentes. \\
    \textbf{C: Comparación}   & Workflows tradicionales, tales como scripts manuales y pipelines estáticos. \\
    \textbf{O: Resultado}      & Nivel de automatización, eficiencia, reproducibilidad y control del flujo de trabajo. \\
    \bottomrule
  \end{tabular}
\end{table}

\subsection{Bases de datos y estrategia de búsqueda}

La búsqueda se realizó en PubMed, complementada con repositorios académicos 
y plataformas de preprints, abarcando publicaciones entre 2020 y 2026. La 
estrategia se ejecutó el 5 de mayo de 2026 mediante la siguiente ecuación:

\begin{quote}\footnotesize\ttfamily\raggedright
("large language models"[TIAB] OR LLM OR "AI agents") \\
\hspace*{1em}AND ("bioinformatics" OR "RNA-seq" OR "genomics") \\
\hspace*{1em}AND ("workflow" OR "pipeline" OR "automation" OR "agent" OR "orchestration")
\end{quote}

Se seleccionaron estudios enfocados en la orquestación de flujos de trabajo y 
sistemas agénticos, integrando adicionalmente documentación técnica sobre 
marcos de trabajo como LangChain y fundamentos de agentes para sustentar el 
análisis arquitectónico.

\subsection{Criterios de inclusión y exclusión}

\paragraph{Criterios de inclusión}

\begin{itemize}
  \item Estudios que describieran arquitecturas explícitas basadas en agentes, incluyendo planificación, uso de herramientas externas o sistemas multi-agente.
  \item Aplicaciones directas en bioinformática, tales como genómica, transcriptómica, proteómica o diseño automatizado de protocolos.
  \item Artículos que analizaran desafíos técnicos asociados al uso de LLMs y agentes, incluyendo alucinaciones, reproducibilidad, coste computacional o control de ejecución.
  \item Publicaciones que propusieran benchmarks, métricas o estrategias de evaluación para sistemas basados en agentes.
\end{itemize}

\paragraph{Criterios de exclusión}

\begin{itemize}
  \item Estudios enfocados exclusivamente en aplicaciones clínicas o diagnóstico por imagen sin relación con workflows bioinformáticos.
  \item Trabajos que utilizaran LLMs únicamente mediante prompting simple, sin capacidades de planificación, uso de herramientas o ejecución autónoma.
  \item Artículos centrados en aplicaciones médicas no relacionadas con bioinformática computacional, incluyendo encuestas de satisfacción de pacientes, medicina tradicional china o estudios clínicos sin componente de automatización de workflows.
  \item Publicaciones sin descripción metodológica clara o sin relación directa con infraestructura computacional basada en agentes.
\end{itemize}

\subsection{Selección de estudios}

La búsqueda inicial recuperó 40 artículos de PubMed. El cribado de títulos,
resúmenes y conclusiones excluyó estudios sobre aplicaciones clínicas,
\textit{prompting} simple o medicina no computacional. Tras aplicar los
criterios de inclusión, se seleccionaron 15 artículos para la síntesis
cualitativa (\cref{fig:prisma}).

Para complementar el análisis de las arquitecturas implementadas en dichos
artículos, se integró de forma dirigida la documentación técnica de LangChain,
la cual sirvió como marco de referencia para fundamentar los componentes de
diseño de agentes, memoria y uso de herramientas discutidos en esta
revisión~\citep{LangChainAgents2026,LangChainMemory2026,LangChainTools2026}.

\begin{figure}[htbp]
  \centering
  \begin{tikzpicture}[
      box/.style  = {rectangle, draw=black!60, fill=blue!5,
                     text width=5cm, minimum height=1cm,
                     align=center, font=\small},
      arrow/.style = {-Stealth, thick, black!60},
      excl/.style  = {rectangle, draw=black!30, fill=red!5,
                      text width=3.5cm, minimum height=0.8cm,
                      align=center, font=\footnotesize}
    ]
    \node[box] (id)     {Registros identificados\\($n = 40$) };
    \node[box, below=0.6cm of id]  (dedup)  {Tras eliminar duplicados\\($n = 40$) };
    \node[box, below=0.6cm of dedup](screen) {Cribado por título/resumen\\($n = 40$) };
    \node[box, below=0.6cm of screen](full)  {Texto completo evaluado\\($n = 15$) };
    \node[box, below=0.6cm of full] (incl)   {Incluidos en la síntesis\\($n = 15$) };

    \node[excl, right=1.2cm of dedup]  (e1) {Duplicados\\($n = 0$) };
    \node[excl, right=1.2cm of screen] (e2) {Excluidos\\($n = 25$) };
    \node[excl, right=1.2cm of full]   (e3) {Excluidos\\($n = 0$) };

    \draw[arrow] (id)     -- (dedup);
    \draw[arrow] (dedup)  -- (screen);
    \draw[arrow] (screen) -- (full);
    \draw[arrow] (full)   -- (incl);
    \draw[arrow] (dedup.east)  -- (e1.west);
    \draw[arrow] (screen.east) -- (e2.west);
    \draw[arrow] (full.east)   -- (e3.west);
  \end{tikzpicture}
  \caption{Diagrama de flujo PRISMA del proceso de selección de estudios.}
  \label{fig:prisma}
\end{figure}
